%!TEX root = ../main.tex
\section{Convex Reformulations}\label{sec:convex-formulations}

Throughout this chapter, we study two-layer, scalar-output ReLU networks
(\cref{eq:two-layer-scalar-output}).
We consider training via regularized ERM as shown in
Problem~\ref{eq:regularized-erm};
for regularization, we use the standard \( \ell_2 \)-squared penalty
\( R(\theta) = \frac{1}{2} \sum_{l=1}^2 \norm{\Wl}_F^2 \), which leads to the
following optimization problem:
\begin{equation}\label{convex-forms:eq:non-convex-relu-mlp}
    \textbf{NC-ReLU}: \, \min_{W^{(1)}, W^{(2)}}   \calL\Big(f_{W^{(1)}, W^{(2)}}(X), y \Big)
    + \frac{\lambda}{2}\sum_{i=1}^m \norm{W^{(1)}_{\cdot, i}}_2^2 + |W^{(2)}_{i}|^2,
\end{equation}
where we recall that \( \lambda \geq 0 \) is the regularization strength.
Problem~\ref{convex-forms:eq:non-convex-relu-mlp} is non-convex, but, as we
have already mentioned, \citet{pilanci2020convexnn} show that there is an
equivalent convex optimization problem with the same optimal value if \( m \geq
m^* \) for some \( m^* \leq n + 1 \).
Furthermore, \citet{wang2022hidden} showed that all optimal solutions to
\eqref{convex-forms:eq:non-convex-relu-mlp} can be found via the convex
problem.

\subsection{Sub-Sampled ReLU Convex Programs}\label{sec:relu-models}

The convex reformulation for the NC-ReLU objective is based on
the possible activations of a single neuron in the hidden layer.
Recall that the activation patterns a ReLU neuron \( \rbr{X w}_{+} \) can take
for fixed \( X \) are described by
\[ \calD_X = \cbr{D = \diag(\mathds{1}(X u \geq 0)) : u \in \R^d}. \]
In the worst case, the number of activation patterns grows exponentially in the
rank of the data matrix \citep{pilanci2020convexnn}.
For any subset \( \tilde \calD \subseteq \calD_X \), we define the sub-sampled
convex optimization problem (C-ReLU):
\begin{equation}\label{convex-forms:eq:convex-relu-mlp}
    \begin{aligned}
        \min_{v, u} \, & \calL\Big(  \sum_{D_i \in \tilde \calD}   D_i X (v_i  -  u_i), y\Big) + \lambda \sum_{D_i \in \tilde \calD} \norm{v_i}_2 + \norm{u_i}_2. \\
                       & \text{s.t.} \quad v_i, u_i \in \calK_i
    \end{aligned}
\end{equation}

\citet{pilanci2020convexnn} prove NC-ReLU and C-ReLU are equivalent using linear semi-infinite duality theory~\citep{ goberna2002semi}.
However, this result requires \( m \geq m^* \) and the full enumeration of the activations of a neuron: \( \tilde \calD = \calD_X \).
In practice, learning with \( \calD_{X} \) is computationally infeasible except for special cases where the data are low rank.
By introducing sub-sampled models, we relax the dependencies on \( m^* \) and \( \calD_X \) to simple inclusions involving \( \tilde \calD \).

%% Moved lemma statement to appendix for space.

%First, we reduce to a group \( \ell_1 \) regularized mixed-integer program using positive homogeneity of the ReLU activation.

%\begin{restatable}{lemma}{miFormulation}\label{lemma:mi-reformulation}
%    The non-convex problem NC-ReLU (Problem \ref{convex-forms:eq:non-convex-relu-mlp}) is equivalent to the mixed-integer program,
%    \begin{equation}\label{convex-forms:eq:mi-reformulation}
%        \begin{aligned}
%            \min_{W^{(1)}, W^{(2)}} \; & \calL\Big(\sum_{i = 1}^m \rbr{X W^{(1)}_{\cdot, i}}_+ W^{(2)}_{i}, y\Big) + \lambda \sum_{i = 1}^m \norm{W^{(1)}_{\cdot, i}}_2 \\
%                               & \text{\emph{s.t.}} \; w_2 \in \cbr{-1, 1}^m.
%        \end{aligned}
%    \end{equation}
%\end{restatable}

%See \cref{app:convex-forms} for a proof.
%\cref{lemma:mi-reformulation} suggests splitting on the values of \( W^{(2)}_{i} \) and then enumerating the ReLU activations functions, which gives the following theorem.

\begin{restatable}{theorem}{convexDualityFree}\label{thm:duality-free}
    Suppose \( \theta^* = \rbr{W^{*(1)}, W^{*(2)}} \) and \( \rbr{v^*, u^*} \)
    are global minima of the NC-ReLU~\eqref{convex-forms:eq:mi-reformulation}
    and C-ReLU~\eqref{convex-forms:eq:convex-relu-mlp} problems, respectively.
    If the number of hidden units satisfies
    \[ m \geq b := \sum_{D_i \in \tilde \calD} \abs{\cbr{v^*_i : v^*_i \neq 0} \cup \cbr{u^*_i : u^*_i \neq 0} }, \]
    and the optimal activations are in the convex model,
    \[ \cbr{\text{\emph{diag}}\rbr{X W^{*(1)}_{\cdot, i} \geq 0 : i \in [m]} } \subseteq \tilde \calD, \]
    then the two problems have same the optimal value.
\end{restatable}
See \cref{app:convex-forms} for proof.
The advantages of this theorem over existing results are (i) the simple and duality-free proof, and (ii) the dependence on \( \tilde \calD \), which we show in \cref{sec:experiments} can be much smaller than \( \calD_X \) while still performing comparably to NC-ReLU.
\cref{thm:duality-free} also reveals that \( m^* \) is determined by the number of active ``neurons'' at the optimal solution of the full C-ReLU problem with \( \tilde \calD = \calD_X \).

%\begin{align*}
%    \rbr{\hat W_1^*, \hat w_2^*}  =   \bigcup_{D_i \in \tilde \calD_i}   \cbr{\rbr{v^*_i, 1} : v^*_i \neq 0} \cup \cbr{\rbr{u^*_i, -1} : u^*_i \neq 0}
%\end{align*}
%is a global minimizer of Problem~\ref{convex-forms:eq:mi-reformulation} and, defining,
%\[
%    T(j) = \cbr{ i \in [m] : \text{diag}\rbr{X W^{(1)}_{\cdot, i}^* > 0} = D_j },
%\]
%we construct a solution to Problem~\ref{convex-forms:eq:convex-relu-mlp} as,
%\begin{equation*}
%    \rbr{\hat v^*, \hat u^*}  =  \bigg\{ \sum_{i \in T(j)}  W^{(1)}_{\cdot, i}^* \mathds{1}(W^{(2)}_{i}^* = 1),    \sum_{i \in T(j)}    W^{(1)}_{\cdot, i}^* \mathds{1}(W^{(2)}_{i}^* = -1) \bigg\}_{j = 1}^m
%\end{equation*}

%\begin{corollary}
%    Combining \cref{thm:duality-free} gives formal conditions for equivalence of C-ReLU and NC-ReLU with \( m \) neurons.
%\end{corollary}

%\begin{remark}\label{remark:active-neurons}
%    Taking \( \tilde \calD = \calD_X \) and applying \cref{thm:duality-free} yields
%    \[ m^* = \sum_{D_i \in \calD_X} \abs{\cbr{v^*_i : v^*_i \neq 0} \cup \cbr{u^*_i : u^*_i \neq 0} }. \]
%\end{remark}

\subsection{Unconstrained Relaxation: Gated ReLUs}

Solving C-ReLU using scalable first-order methods typically requires projecting on \( \calK_i \), which is an expensive quadratic program in the general case.
To circumvent this, we consider the following unconstrained relaxation:
\begin{equation}\label{convex-forms:eq:unconstrained-relu-mlp}
    \textbf{C-GReLU}: \, \min_{w} \; \calL\Big(\sum_{D_i \in \tilde \calD} D_i X w_i, y\Big)+ \lambda \sum_{D_i \in \tilde \calD} \norm{w_i}_2.
\end{equation}
At first look, this problem is a high-dimensional GLM with group-\( \ell_1 \) regularization.
In fact, C-GReLU is the convex re-formulation of another neural network optimization problem.
Let \( \calG \subset \R^d \) and consider the model,
\[
    f_{W^{(1)}, W^{(2)}}(X) = \sum_{g_i \in \calG} \phi_{g}(X, W^{(1)}_{\cdot, i}) W^{(2)}_{i},
\]
where \( \phi_{g}(X, u) = \text{diag}(\mathds{1}(Xg \geq 0)) X u \) is a ``gated ReLU'' activation function with fixed gate vector \( g \) \citep{fiat2019decoupling}.
C-GReLU is equivalent to training this gated ReLU network.

\begin{restatable}{theorem}{unconstrainedEquivalence}\label{thm:unconstrained-equivalence}
    Let \( g_i \in R^d \) such that \(\text{diag}(X g_i \geq 0) = D_i \) and \(
    \tilde \calG = \cbr{g_i : D_i \in \tilde \calD} \).
    Then, C-GReLU is equivalent to the following gated ReLU problem:
    \begin{equation}\label{convex-forms:eq:gated-relu-mlp}
        \textbf{NC-GReLU}: \, \min_{W^{(1)}, W^{(2)}} \calL\Big(\sum_{g_i \in \tilde \calG}
        \phi_{g_i}(X, W^{(1)}_{\cdot, i})W^{(2)}_{i}, y\Big) + \frac{\lambda}{2}
        \sum_{g_i \in \tilde \calG} \norm{W^{(1)}_{\cdot, i}}_2^2 + \abs{W^{(2)}_{i}}^2.
    \end{equation}
\end{restatable}
See \cref{app:convex-forms} for proof.
We can use \cref{thm:unconstrained-equivalence} to fit Gated ReLU networks by (much easier) unconstrained minimization.
However, as the next section shows, we can also leverage C-GReLU to approximate or exactly solve the original ReLU problem.
